\documentclass[../../../../main.tex]{subfiles}
\begin{document}

\begin{flushright}
\footnotesize\textit{【自動文字起こし・要確認】}
\end{flushright}

[解] 3交点の $x$ 座標を、小さい順に $\alpha, \beta, \gamma$ とおく。\\
又、$f(x) = ax^2 + bx + c$ とする。
\[
x^3 - f(x) = (x-\alpha)(x-\beta)(x-\gamma) \cdots \text{\textcircled{1}}
\]
各点での接線を $l_1, l_2, l_3$ として、これらは
\[
y = 3k^2 x - 2k^3 \quad (k=\alpha, \beta, \gamma) \cdots \text{\textcircled{2}}
\]
で表される。\textcircled{2} が $P$ を通ることから
\[
2k^3 - 3pk^2 + q = 0 \cdots \text{\textcircled{3}}
\]
したがって、$\alpha, \beta, \gamma$ は $k$ の3次式 \textcircled{3} の3実解で、\textcircled{1} から
\[
x^3 - f(x) = x^3 - \frac{3}{2}p x^2 + \frac{q}{2}
\]
係数比較して
\[
a = \frac{3}{2}p, \quad b = 0, \quad c = -\frac{q}{2} \cdots (1)
\]

(2) $p, q$ の条件は $x$ の3次方程式 $x^3 - \frac{3}{2}px^2 + \frac{q}{2} = 0 \cdots \text{\textcircled{4}}$ が3異実解を持つこと。\textcircled{4} の左辺 $g(x)$ とおく。$g'(x) = 3x^2 - 3px$ から、$p$ によって下のように分かれる

$1^\circ \ p = 0$ の時\\
$g'(x) \ge 0$ となり、$g(x)$ は単調増加し、\textcircled{4} が3異実解を持つことはない。

$2^\circ \ p > 0$ の時\\
\begin{center}
\begin{tabular}{|c|c|c|c|c|c|}
\hline
$x$ & $\cdots$ & $0$ & $\cdots$ & $p$ & $\cdots$ \\
\hline
$g'$ & $+$ & $0$ & $-$ & $0$ & $+$ \\
\hline
$g$ & $\nearrow$ & & $\searrow$ & & $\nearrow$ \\
\hline
\end{tabular}
\end{center}
上図から、条件は
\begin{align*}
& g(0) > 0 \land g(p) < 0 \\
\iff & \frac{q}{2} > 0 \land - \frac{p^3}{2} + \frac{q}{2} < 0 \\
\iff & q > 0 \land -p^3 + q < 0
\end{align*}

$3^\circ \ p < 0$ の時\\
$2^\circ$ と同様に、条件は
\[
g(p) > 0 \land g(0) < 0 \iff q < 0 \land -p^3 + q > 0
\]

以上まとめて
\[
\begin{cases}
p > 0, \quad 0 < q < p^3 \\
\text{又は} \\
p < 0, \quad p^3 < q < 0
\end{cases}
\]

\end{document}
